% Part: first-order-logic
% Chapter: proof-systems
% Section: seqeunt-calculus

\documentclass[../../../include/open-logic-section]{subfiles}

\begin{document}

\iftag{FOL}
      {\olfileid{fol}{prf}{seq}}
      {\olfileid{pl}{prf}{seq}}

\olsection{De sequentencalculus}

Veel !!{derivation}ssystemen werken met rangschikkingen van
!!{sentence}s; de sequentencalculus werkt daarentegen met
\emph{sequenten}. Een sequent is een uitdrukking van de vorm
\[
!A_1, \dots, !A_m \Sequent !B_1, \dots, !B_m,
\]
dat wil zeggen: een paar rijen !!{sentence}s, van elkaar gescheiden
door het sequententeken~$\Sequent$. Elk van beide rijen mag leeg zijn.
!!^a{derivation} in de sequentencalculus is een boom van sequenten. De
bovenste sequenten hebben een bijzondere vorm (ze heten
``beginsequenten'' of ``axioma's''), en elke andere sequent volgt
volgens een van de afleidingsregels uit de sequenten die er direct
boven staan. De afleidingsregels veranderen de !!{sentence}s in de
sequenten (door ze links of rechts toe te voegen, weg te laten of te
verplaatsen), of voeren in de conclusie van de regel een samengestelde
!!{formula} in. Zo staat de regel $\LeftR{\land}$ de afleiding van $!A,
\Gamma \Sequent \Delta$ naar $!A \land !B, \Gamma \Sequent \Delta$ toe,
en $\RightR{\lif}$ de afleiding van $!A, \Gamma \Sequent \Delta, !B$
naar $\Gamma \Sequent \Delta, !A \lif !B$, voor alle $\Gamma$, $\Delta$,
$!A$ en~$!B$. (In het bijzonder mogen $\Gamma$ en~$\Delta$ leeg zijn.)

De op de sequentencalculus gebaseerde relatie $\Proves$ wordt als
volgt gedefinieerd: $\Gamma \Proves !A$ dan en slechts dan als er een
rij $\Gamma_0$ bestaat zodat elke $!A$ in~$\Gamma_0$ tot~$\Gamma$
behoort en er een
!!{derivation} bestaat met de sequent~$\Gamma_0 \Sequent !A$ aan de
wortel. $!A$ is een stelling in de sequentencalculus als de
sequent~$\Sequent !A$ !!a{derivation} heeft. De volgende
!!a{derivation} laat bijvoorbeeld zien dat $\Proves (!A \land !B) \lif
!A$:
\begin{prooftree}
  \Axiom$!A \fCenter !A$
  \RightLabel{\LeftR{\land}}
  \UnaryInf$!A \land !B \fCenter !A$
  \RightLabel{\RightR{\lif}}
  \UnaryInf$\fCenter (!A \land !B) \lif !A$
\end{prooftree}

Een verzameling $\Gamma$ is inconsistent in de sequentencalculus als
er !!a{derivation} van $\Gamma_0 \Sequent$ bestaat (waarbij elke $!A
\in \Gamma_0$ tot~$\Gamma$ behoort en de rechterkant van de sequent
leeg is). Met de regel \RightR{\Weakening} kan uit een inconsistente
verzameling elke !!{sentence} worden !!{derive}d.

Gerhard Gentzen vond de sequentencalculus in de jaren dertig van de
twintigste eeuw uit. Door zijn systematische en symmetrische opzet is
hij zeer geschikt om een theorie van !!{derivation}s te ontwikkelen.
In de sequentencalculus zijn !!{derivation}s betrekkelijk gemakkelijk
te vinden, maar ze zijn vaak moeilijk te lezen en hun verband met
bewijzen is niet altijd gemakkelijk te doorzien. Niettemin is de
sequentencalculus een zeer elegante benadering van
!!{derivation}ssystemen gebleken en bestaan er voor veel logica's
sequentencalculi.
\end{document}
